\documentclass[11pt]{article}

\usepackage[margin=1in]{geometry}
\usepackage[utf8]{inputenc}
\usepackage{hyperref}
\usepackage{graphicx}
\usepackage{booktabs}
\usepackage{amsmath}
\usepackage{amssymb}
\usepackage{xcolor}
\usepackage{natbib}

\hypersetup{
  colorlinks=true,
  linkcolor=blue,
  citecolor=blue,
  urlcolor=blue
}

\title{The StackGraph: A Persistent Neuro-Symbolic World-Model for\\
Agentic Root-Cause Analysis at Homelab Cost}
\author{Victor Robin\\ BlueRobin (independent homelab research)}
\date{June 2026}

\begin{document}

\maketitle

\begin{abstract}
We describe the StackGraph, a single persistent property graph that fuses five
disparate layers of a self-hosted microservice platform---static code topology,
static infrastructure topology, runtime telemetry service-graph, change/deploy
history, and accumulated root-cause-analysis (RCA) outcomes including a
bi-temporal incident memory---into one world-model stored as a second graph key
inside an existing FalkorDB instance. The StackGraph is consulted by an
otherwise graph-free, ReAct-style large-language-model (LLM) agent
\citep{yao2022react} that performs advisory RCA on production incidents in the
BlueRobin homelab. The design occupies a deliberate neuro-symbolic middle
ground: deterministic graph traversal proposes ranked candidate root causes, and
the LLM narrates and verifies them against real tool evidence; it never grades
its own homework. Every node and edge carries a first-class provenance bundle
(source, status, confidence) so that any ranked candidate is traceable to the
telemetry, code, or change that produced it. We report the schema, the $k$-hop
neighbourhood read path, the parameterised-Cypher security posture, and the
engineering discipline that keeps the entire stack---graph, vector index, and
temporal memory---inside a $\sim$50~EUR/month budget with no GraphRAG community
indexing and no new always-on workload. This is a systems and experience report,
not an empirical breakthrough: the system is advisory and human-in-the-loop by
construction, and the internal evaluation is an admittedly oversimplified
self-seeded benchmark (companion paper). The contribution is the architecture
and the cost discipline.
\end{abstract}

\section{Introduction}

Automated root-cause analysis for microservice and Kubernetes systems is hard,
and recent independent benchmarks make the difficulty concrete rather than
rhetorical. The best agent on OpenRCA reaches roughly 11.3\% exact-match
\citep{xu2025openrca}, and the best frontier model on IBM's ITBench-AA reaches
roughly 47\% precision-at-full-recall \citep{itbenchaa2026}. These ceilings are
the honest baseline against which any homelab system must be read: fully
autonomous RCA is not a solved problem, and an advisory, human-in-the-loop
posture is the state of the art, not a concession.

A naive response is to hand a frontier LLM the alert text and a pile of tools
and let it reason. This fails in well-documented ways. The MAST failure taxonomy
\citep{cemri2025mast}, derived from 310 ITBench traces, attributes roughly 94\%
of multi-agent failures to a handful of modes---incorrect self-verification,
premature termination, memory loss, and reasoning-action mismatch. An LLM asked
to localise a fault from symptoms alone tends to hallucinate a plausible-sounding
cause, declare victory, and terminate. The RCA survey literature reaches the same
conclusion structurally: LLM-only RCA ``lacks structural grounding and risks
hallucinated or unsafe suggestions'' \citep{zhang2024survey}. The fix the
literature converges on is grounding the agent in an explicit topology or causal
graph.

But graph grounding has a cost profile that does not fit a hobbyist budget.
Microsoft GraphRAG \citep{edge2024graphrag} delivers strong global sensemaking
but is dominated by LLM indexing calls---entity, edge, and claim extraction plus
community summarisation---at roughly 20--40~USD per million tokens, re-run
whenever the corpus changes. A persistent agent-memory graph such as Zep/Graphiti
\citep{rasmussen2025zep} is the right pattern but ships as a Python/TypeScript
sidecar. Neither fits a single memory-saturated node running a .NET service under
a hard 20~EUR/month LLM ceiling.

The StackGraph is our answer: a persistent, graph-grounded world-model built and
queried entirely deterministically, with the LLM demoted to a narrator and
verifier over graph-proposed candidates. Its contributions are:

\begin{itemize}
  \item \textbf{A five-layer fused world-model in one property graph.} Code
  topology, infra topology, runtime telemetry, change history, and RCA outcomes
  (including bi-temporal memory) live in a single FalkorDB graph keyed
  \texttt{bluerobin\_stack}, not five siloed stores.
  \item \textbf{A neuro-symbolic read path.} Deterministic $k$-hop traversal and
  personalised-PageRank blame-propagation propose ranked candidates that seed a
  graph-free-style ReAct loop; the LLM confirms leads against hard tool evidence
  and never validates a hypothesis on its own inference.
  \item \textbf{First-class provenance on every edge.} A
  \texttt{(source, status, confidence)} bundle (ADR-021) makes every candidate
  auditable back to the telemetry, code unit, or deploy that produced it.
  \item \textbf{A homelab cost envelope.} A second graph key in the
  already-deployed FalkorDB, local Ollama embeddings that bypass the paid edge,
  no GraphRAG community indexing, and no new always-on workload---the whole
  world-model is cost-neutral within $\sim$50~EUR/month.
  \item \textbf{A fail-open advisory contract.} An empty, stale, or unreachable
  graph degrades to an empty plan and the agent proceeds on its tool floor
  (NFR-SG-2); the graph is never a hard dependency of incident response.
\end{itemize}

\section{Related Work}

\paragraph{Two architectural camps.}
Contemporary RCA systems split into graph-grounded ``system-model'' builders and
graph-free live investigators. The grounded camp includes Dynatrace Davis
(deterministic causal traversal of the Smartscape graph), Causely (a Bayesian
causal DAG plus codebook), Traversal (a learned causal dependency map),
Resolve.ai (a 50k-node dynamic knowledge graph with PR citations), Deductive AI
(a Neo4j code-to-telemetry graph with reinforcement learning), and Cleric (a
hierarchical knowledge graph driving massively parallel multi-hypothesis search).
The graph-free camp---HolmesGPT (Robusta and Microsoft, CNCF Sandbox), K8sGPT,
Datadog Bits AI SRE, Grafana Assistant, Parity, and NeuBird Hawkeye---correlates
at investigation time with no persistent model. The StackGraph deliberately
occupies the middle ground identified in the survey framing: graph traversal
proposes candidates and the LLM narrates and verifies, rather than the LLM being
either the sole reasoner or a thin shell over a symbolic engine.

\paragraph{GraphRAG and temporal knowledge graphs.}
Microsoft GraphRAG \citep{edge2024graphrag} introduced hierarchical Leiden
community detection with LLM-generated community summaries for global sensemaking,
at an indexing cost dominated by LLM calls. Zep/Graphiti
\citep{rasmussen2025zep} is the pattern we follow for memory: bi-temporal edges
carrying both valid time and ingestion time, with hybrid retrieval combining
semantic cosine similarity, BM25, and breadth-first graph traversal and crucially
no LLM at retrieval (Graphiti reports P95 around 300~ms). We adopt the
bi-temporal schema and the no-LLM-at-retrieval discipline but hand-roll them in
.NET on FalkorDB rather than deploying the Python \texttt{graphiti-core} library.
For RCA specifically, SynergyRCA \citep{liu2025synergyrca} pairs a
spatio-temporal StateGraph with a MetaGraph in Neo4j/Cypher, retrieval-augmented
to an LLM, reaching roughly 0.90 precision in around two minutes on two
production clusters---the strongest credible academic KG-plus-LLM RCA result, and
a useful upper reference for what graph grounding buys. The serialization of
graph context for the LLM follows ``Simple Is Effective'' \citep{li2024simple},
which shows that rendering triples as template natural-language sentences (``the
relation of head is tail'') improves LLM comprehension over raw triples.

\paragraph{Grounding versus hallucination.}
The RCA survey of \citet{zhang2024survey} and the MAST taxonomy
\citep{cemri2025mast} together motivate the verification discipline: never let
the LLM grade its own homework, require hard tool evidence before a confident
exit, and enforce termination via a finite-state machine outside the model. The
StackGraph supplies the symbolic substrate those disciplines run against; the
verification gate and hypothesis machinery that consume it are detailed in the
companion papers in this series.

\section{Method: The StackGraph World-Model}

\subsection{What the StackGraph is}

The StackGraph is not a stack-trace graph or a function-level call graph. It is a
persistent topology-and-outcome world-model of the entire platform, keyed
\texttt{bluerobin\_stack} in FalkorDB (FR-SG-1), that fuses five layers:

\begin{enumerate}
  \item \textbf{Static code topology}---services, FastEndpoints endpoints, NATS
  subjects, and structure-defining source files.
  \item \textbf{Static infra topology}---namespaces, deployments, nodes, Traefik
  routes, datastores, and external APIs.
  \item \textbf{Runtime telemetry}---the Tempo service-graph projected as
  \texttt{CALLS} and \texttt{READS\_FROM} edges carrying error-rate and latency.
  \item \textbf{Change history}---\texttt{Deploy} and \texttt{Change} nodes
  derived from Flux reconciliations and Git commits.
  \item \textbf{RCA outcomes}---\texttt{Incident}, \texttt{Alert},
  \texttt{Hypothesis} nodes, evidence and root-cause edges, plus a bi-temporal
  incident-memory subgraph.
\end{enumerate}

Conceptually (Figure described in prose): an Alertmanager alert and the
\texttt{bluerobin\_stack} graph both feed the \texttt{StackGraphRcaService}
$k$-hop read path. The five layers live inside the one FalkorDB property graph.
The read path emits a ranked candidate plan---paths plus evidence plus suggested
tools---which seeds a graph-free-style ReAct loop in which the LLM narrates and
verifies.

\subsection{Schema and provenance}

Every node and edge carries an ADR-021 provenance bundle as first-class metadata,
expressed by three enums:

\begin{verbatim}
public enum EdgeSource { StaticCode, StaticInfra, Tempo, Loki,
                         Prometheus, Synthetic, LlmInferred, Git }
public enum EdgeStatus { Active, Stale, Candidate, Deprecated }
public enum RiskClass  { Metadata, Structural }
\end{verbatim}

Standard properties seen across builders are \texttt{source}, \texttt{status},
\texttt{confidence\_score}, \texttt{first\_seen} (set \texttt{ON CREATE}),
\texttt{last\_seen}, and \texttt{last\_scan}. The node labels and their natural
keys:

\begin{table}[h]
\centering
\small
\begin{tabular}{lll}
\toprule
\textbf{Label} & \textbf{Natural key} & \textbf{Notes} \\
\midrule
\texttt{Service} & \texttt{name} & criticality, kind (infra/worker), namespace, repo \\
\texttt{Endpoint} & \texttt{service|method|route} & FastEndpoints route \\
\texttt{NatsSubject} & \texttt{subject} & event-bus subject \\
\texttt{CodeUnit} & \texttt{repo|path} & \texttt{embedding\_ref} = Qdrant point id, \texttt{github\_url}, \texttt{language} \\
\texttt{Datastore} & \texttt{name} & \texttt{db\_system} \\
\texttt{ExternalApi} & \texttt{name} & GitHub / R2 / CF AI Gateway peers \\
\texttt{Deployment} & \texttt{name} & namespace, image, replicas \\
\texttt{Node} & \texttt{name} & cluster node, role \\
\texttt{Namespace} & \texttt{name} & environment \\
\texttt{Route} & \texttt{host} & Traefik \\
\texttt{Alert} & \texttt{fingerprint} & Alertmanager fingerprint \\
\texttt{Incident} & \texttt{fingerprint} & RCA write-back; \texttt{evidence\_chain\_json} \\
\texttt{Deploy} & \texttt{service|image\_digest|deployed\_at} & from Flux \\
\texttt{Change} & \texttt{repo|commit\_sha} & from Git; \texttt{pr\_number}, \texttt{summary} \\
\texttt{Hypothesis} & \texttt{incident\_fingerprint|suspect\_ref} & \texttt{state}, \texttt{rank} \\
\texttt{MemoryEpisode} & \texttt{incident\_fingerprint} & bi-temporal memory \\
\texttt{MemoryEntity} & \texttt{entity\_type|entity\_ref} & deduped fact entity \\
\bottomrule
\end{tabular}
\caption{StackGraph node labels and natural keys.}
\end{table}

Edge types span the five layers: \texttt{EXPOSES},
\texttt{PUBLISHES\_TO}/\texttt{CONSUMES\_FROM}/\texttt{DECLARES\_SUBJECT} (NATS),
\texttt{CALLS} (carrying \texttt{trace\_count}/\texttt{error\_rate}/%
\texttt{latency\_p95\_ms}), \texttt{READS\_FROM}, \texttt{IMPLEMENTS},
\texttt{DEFINED\_IN}, \texttt{DEPLOYED\_AS}/\texttt{RUNS\_ON}/%
\texttt{IN\_NAMESPACE}/\texttt{ROUTES\_TO}, \texttt{DEPLOYED}/\texttt{CHANGED},
\texttt{FIRED\_ON}/\texttt{AFFECTS}, \texttt{EVALUATED\_AS},
\texttt{EVIDENCED\_BY} (carrying \texttt{tool}/\texttt{verdict}/%
\texttt{result\_summary}), \texttt{ROOT\_CAUSED\_BY} (written only when a
hypothesis is validated and confidence $\geq 0.40$), \texttt{CITES} (Hypothesis
to Change/Deploy), and the bi-temporal memory edges
\texttt{MENTIONS}/\texttt{RESOLVED\_BY}/\texttt{SUPERSEDES} carrying
\texttt{t\_valid}/\texttt{t\_invalid}/\texttt{ingested\_at}.

The design principle behind the bundle is auditability: because \texttt{source}
distinguishes \texttt{StaticInfra} from \texttt{Tempo} from \texttt{LlmInferred},
a ranked candidate can always be traced to the evidence class that produced it,
and LLM-inferred structure can be structurally barred from validating an outcome.

\subsection{How the graph is built: a single writer}

The graph has exactly one writer, \texttt{GraphIngestor}, implemented as a
partial class split across ingest lanes. Its contract is idempotent
MERGE-on-natural-key, which makes every rebuild-from-source deterministic and
re-runnable. The lanes are:

\begin{itemize}
  \item \textbf{Static code lane} clones each mapped repo and parses FastEndpoints
  routes and NATS publish/subscribe symbols. Only structure-defining files become
  \texttt{:CodeUnit} nodes (not every \texttt{.cs} file) so the graph stays
  readable, linking \texttt{CodeUnit-[:IMPLEMENTS]->Service},
  \texttt{Endpoint-[:DEFINED\_IN]->CodeUnit}, and
  \texttt{CodeUnit-[:DECLARES\_SUBJECT]->NatsSubject} (FR-SG-2).
  \item \textbf{Static infra lane} reads Kubernetes and Traefik CRDs into
  Namespace/Deployment/Node/Route/Datastore nodes.
  \item \textbf{Telemetry lane} projects the Tempo service-graph metrics into
  \texttt{CALLS}/\texttt{READS\_FROM} edges with the runtime bundle, and flips an
  edge to \texttt{status = Stale} once its metric ages past a freshness threshold
  (FR-SG-3, $\leq 15$~min).
  \item \textbf{Change lane} writes \texttt{Deploy} nodes from Flux and
  \texttt{Change} nodes from Git.
  \item \textbf{RCA write-back} records the concluded \texttt{Incident}, an
  \texttt{AFFECTS} edge to the symptom service always, and a
  \texttt{ROOT\_CAUSED\_BY} edge to the cause only at confidence $\geq 0.40$
  (FR-SG-7).
\end{itemize}

The single-writer-many-readers model (ADR-021) avoids write contention and keeps
the rebuild path a pure function of its sources.

\subsection{The $k$-hop read path}

RCA reads the graph through \texttt{StackGraphRcaService.InvestigateAsync}
(FR-SG-6). Given an alerting service it pulls the $k$-hop neighbourhood
(parameterised, SEC-SG-4), turns each candidate into a topology/telemetry/history
scoring signal, and returns a ranked top-$N$ investigation plan of paths,
evidence, and suggested tools. The simplest collection step is a one-hop Cypher
query over the runtime call edges:

\begin{verbatim}
MATCH (s:Service {name: $name})-[e:CALLS]->(n:Service)
RETURN n.name, n.criticality, e.error_rate, e.status
\end{verbatim}

Each neighbour becomes a \texttt{ScoringSignal} whose topology weight tracks
criticality, whose telemetry weight tracks error-rate (halved when the edge is
\texttt{Stale}), and whose history weight comes from a root-cause-history prior,
together with a \texttt{Path} and a set of \texttt{SuggestedTools}. A richer
bounded multi-hop traversal exists for alert-anchored search:

\begin{verbatim}
MATCH (a:Alert {fingerprint: $fingerprint})-[:AFFECTS]->(svc:Service)
MATCH path = (svc)-[*1..4]->(candidate:Service)
RETURN candidate, path LIMIT $bounded
\end{verbatim}

The hop count is clamped at four. Candidates are ranked by a deterministic
blame-propagation scorer---personalised PageRank seeded at the symptom over the
error-weighted dependency subgraph---and a multi-term blend, both detailed in the
companion ranking paper. The output is strictly advisory: an empty, stale, or
unreachable graph yields an empty plan (NFR-SG-2), and the agent proceeds on its
tool floor.

\subsection{The neuro-symbolic contract}

The candidates from the read path are not conclusions; they are leads. They seed
the LLM prompt verbatim under a header that frames them as ``likely root causes
(graph blame-propagation, before tools)---prioritised leads to confirm, not
conclusions.'' The LLM then runs a classic ReAct tool loop
\citep{yao2022react} over real observability tools (Tempo, Loki, Prometheus,
Kubernetes, Flux, NATS, GitHub, and a \texttt{query\_stack\_graph} tool), and a
hypothesis's state is derived from the tool evidence it accumulates, never
asserted by the model. Evidence whose provenance \texttt{source} is
\texttt{LlmInferred} is rejected at the sink---the system never grades its own
homework. This is the operational meaning of the neuro-symbolic middle ground:
symbolic traversal narrows the search space cheaply and deterministically; the
neural component narrates, investigates, and verifies, but cannot self-certify.
The serialization of a candidate path into template natural-language sentences
for the LLM follows ``Simple Is Effective'' \citep{li2024simple} and is exposed,
A/B-gated, through the \texttt{query\_stack\_graph} tool.

\section{Implementation}

\subsection{Storage and query}

The StackGraph lives in FalkorDB (\texttt{falkordb/falkordb:v4.4.1}, 2~GiB limit,
Redis wire protocol on port 6379) as a second graph key alongside the
application's existing graph (ADR-021, FR-SG-1, NFR-SG-1). The .NET client drives
it over the RedisGraph wire protocol via \texttt{GRAPH.QUERY} using
\texttt{StackExchange.Redis}:

\begin{verbatim}
return db.Execute("GRAPH.QUERY", graphKey, composed, "--compact");
\end{verbatim}

All queries are constant parameterised Cypher templates. Parameters bind through
FalkorDB's \texttt{CYPHER name=value} preamble---never string interpolation---%
which is the FalkorDB analogue of parameterised SQL (SEC-SG-4). Values are
emitted as escaped Cypher literals so that an injection payload becomes the
literal data of a single string parameter, and the relationship types and labels
that Cypher forbids parameterising are always selected from fixed code-supplied
allow-lists (for example, suspect labels collapse to a seven-label set) rather
than from caller text. This closes both the Cypher-injection and the
open-redirect/SSRF surfaces that a caller-influenced graph query would otherwise
open.

\subsection{Embeddings and the code-to-graph join}

The vector store is a separate Qdrant instance reached over gRPC (port 6334,
cosine distance). Embeddings are generated against an in-cluster Ollama
OpenAI-compatible \texttt{/v1/embeddings} endpoint, which is direct and free and
never transits the paid Cloudflare edge (ADR-020). Two pipelines coexist: code
retrieval uses \texttt{qwen3-embedding:8b} (4096-d) into a
\texttt{code-embeddings} collection, and incident memory uses \texttt{bge-m3}
(1024-d) into a dedicated collection (ADR-047). The static code lane embeds each
\texttt{:CodeUnit} and persists the returned Qdrant point id onto
\texttt{CodeUnit.embedding\_ref} (FR-SG-48), enabling bidirectional traversal:
graph-to-code returns the embedding refs for a service's code units, and
code-to-graph maps a semantic code hit back to its owning \texttt{:Service}.
Chunking is Roslyn-based, by method, constructor, class, record, interface, enum,
and computed property.

\subsection{The $\sim$50~EUR/month cost envelope}

The cost discipline is the load-bearing engineering claim, and it is enforced by
what the design refuses to do. The platform runs to a hard total ceiling of
50~EUR/month: a single Hetzner CX43 worker at roughly 17.5~EUR, Cloudflare R2 and
Tunnel on free tiers, and a hard 20~EUR/month LLM token ceiling (COST-7) with a
50/80/100\% alert ladder that fail-closes new analyses at 100\%. The hosting node
is memory-saturated (around 96\% of memory requests committed), which is the
binding reason the StackGraph must live as a second key inside the existing 2~GiB
FalkorDB rather than as a new pod (NFR-SG-1). The architecture consequently
rejects, on cost grounds, four otherwise-attractive options: GraphRAG community
indexing \citep{edge2024graphrag} (its 20--40~USD per million indexing tokens
would breach the LLM ceiling on its own), a Python \texttt{graphiti-core} sidecar
\citep{rasmussen2025zep} (a new always-on workload the node cannot host),
foundation-model anomaly detection (10--50$\times$ slower and no accuracy win),
and any second time-series database (the graph indexes over the existing LGTM
stack rather than re-storing telemetry). The result is a graph-grounded,
temporally-aware RCA world-model whose marginal monthly cost is effectively
zero---it adds a graph key, a Qdrant collection, and a periodic CronJob ingest
(NFR-SG-3), and nothing that runs continuously.

\subsection{Operational posture}

The read path is health-gated and fail-open (NFR-SG-2): the agentic loop has an
eleven-tool floor and a five-minute run budget, and the StackGraph is one input
among them, never a precondition. Ingest runs as a cost-neutral CronJob
(NFR-SG-3). The write-back of incident outcomes and bi-temporal memory episodes
is best-effort and shadow-first, so the whole world-model can be populated and
observed in production before any of its rankings are allowed to govern a
decision. All of this is governed by FEAT-027 (Stack-Graph RCA), with the graph
model and LLM-refinement safety captured in ADR-021 and ADR-022 respectively.

\section{Evaluation}

This paper's contribution is architectural; the quantitative evaluation belongs
to the companion eval paper, and we summarise it here only with its caveats
stated plainly, because honest framing matters more than a headline number.

The internal harness replays a self-seeded set of nine labelled incidents through
the real RCA read path and emits a deterministic scorecard. On that set the
system reports accuracy-at-1 (AC@1, the fraction of incidents whose top-ranked
candidate is the labelled cause) of 1.0, and mean-time-to-resolution (MTTR)
figures in the sub-second range.

\begin{table}[h]
\centering
\small
\begin{tabular}{lll}
\toprule
\textbf{Metric} & \textbf{Value} & \textbf{What it actually is} \\
\midrule
AC@1 (self-seeded set) & 1.0 & Top-ranked candidate matches the label on 9 hand-labelled incidents \\
Incident count & 9 & Authors' own set; described internally as ``oversimplified'' \\
MTTR proxy & sub-second & Replay-clock wall time, not production incident MTTR \\
External ref.: OpenRCA best agent & $\sim$11.3\% & Exact-match, 335 cases \citep{xu2025openrca} \\
External ref.: ITBench-AA best model & $\sim$47\% & Precision@full-recall, 59 tasks \citep{itbenchaa2026} \\
\bottomrule
\end{tabular}
\caption{Internal scorecard with honest external reference ceilings.}
\end{table}

These numbers must not be read as a capability claim. AC@1 = 1.0 on nine
self-labelled incidents is a regression sentinel, not evidence of general RCA
accuracy; the authors themselves characterise the set as an oversimplified
benchmark. The MTTR figures are replay-clock proxies measured against synthetic
telemetry, not production mean-time-to-resolution. The honest comparison points
are the external ceilings in the table: against OpenRCA's $\sim$11\%
\citep{xu2025openrca} and ITBench-AA's $\sim$47\% \citep{itbenchaa2026}, an
advisory homelab system is matching the posture of the field, not exceeding it. A
distractor-injection arm of the harness (candidates with higher topology
criticality but lower error-rate) exists precisely so that AC@1 below 1.0 is
achievable and the gate can catch real regressions rather than rubber-stamping a
trivially separable set. The defensible claim is narrow and architectural: the
StackGraph read path produces a ranked, provenance-tagged plan deterministically
and cheaply, and it does so without ever letting the LLM certify its own
conclusion.

\section{Limitations and Threats to Validity}

\paragraph{Construct validity.}
The evaluation set is self-seeded and small (nine incidents), authored by the
same people who built the system, and acknowledged internally as oversimplified.
AC@1 = 1.0 on it says little about real incident diversity. The MTTR figures are
replay-clock proxies, not production MTTR, and should never be quoted as
operational latency.

\paragraph{Internal validity.}
Because ingest is deterministic and self-seeded, the labelled incidents and the
graph that ranks them share provenance; a separable-by-construction set inflates
AC@1. The distractor arm mitigates but does not eliminate this. The system is
advisory and human-in-the-loop by design---a human always merges---so end-to-end
``resolution'' is never autonomous and the harness measures localisation, not
repair.

\paragraph{External validity.}
The platform is a single-node homelab with a small service count; the $k$-hop
neighbourhood is small and the four-hop clamp has never been stressed at scale.
Findings may not transfer to large multi-cluster estates where graph size,
community structure, and fan-out dominate---exactly the regime where
GraphRAG-style indexing \citep{edge2024graphrag} earns its cost.

\paragraph{Known engineering caveats.}
Two are worth flagging. First, the create-on-demand path for the incident-memory
Qdrant collection can inherit the code-index embedding dimension (4096-d) while
the documented memory model is \texttt{bge-m3} at 1024-d, a latent mismatch
unless the declarative collection-init job wins the race; the dimensions must be
pinned (the intent of ADR-047). Second, the graph is only as fresh as its ingest
CronJob and its staleness thresholds; a stale \texttt{CALLS} edge is
down-weighted, not removed, so a candidate can persist on topology weight after
its telemetry has gone cold. These are mitigations, not guarantees, and are
stated so that the world-model is not over-trusted.

\section{Conclusion}

The StackGraph shows that graph-grounded agentic RCA does not require a
graph-indexing budget. By fusing five topology and outcome layers into a single
FalkorDB property graph keyed \texttt{bluerobin\_stack}, carrying a provenance
bundle on every edge, and consulting it through a deterministic $k$-hop read path
that seeds a graph-free-style ReAct loop \citep{yao2022react}, the system gets
the structural grounding the literature calls for \citep{zhang2024survey} while
demoting the LLM to a narrator and verifier that cannot grade its own homework.
The whole world-model---graph, vector join, and bi-temporal memory---fits inside
an existing FalkorDB instance on a memory-saturated single node under a
50~EUR/month ceiling, with no community indexing, no Python sidecar, and no new
always-on workload. The honest scope is equally important: the system is advisory
by construction and its internal evaluation is a small self-seeded sentinel, not
a capability proof. The neuro-symbolic middle ground it occupies---symbolic
traversal proposes, neural reasoning verifies---is, we argue, the right shape for
RCA at the autonomous ceiling the field actually sits at today
\citep{xu2025openrca,itbenchaa2026}. Subsequent papers in this series detail the
correlation-first ranking blend, the externalised verification gate and
hypothesis tree, the eval-first methodology, and the bi-temporal incident memory
\citep{rasmussen2025zep} that the StackGraph makes possible.

\bibliographystyle{plainnat}
\bibliography{references}

\end{document}
